\section{Completion and quotient repair}\label{sec:repair}

Initial semantics asks what every model must satisfy.  Carrier-valued completion asks for the existence of a much more special model, one in which no named action value leaves the carrier.  The distinction is already visible in one protected row.

\subsection{Completion is splitting}

Let
\[
\begin{array}{ccc}
S & \hookrightarrow & Q_r\\
{\scriptstyle h}\downarrow && \downarrow{\scriptstyle p_r}\\
Q_* & \xrightarrow{m} & P
\end{array}
\]
be a protected row pushout, so $m$ is injective.

\begin{theorem}[Completion as a split monomorphism]\label{thm:splitting}
The partial homomorphism $h:S\to Q$ extends to an endomorphism $\tau:Q\to Q$ if and only if the distinguished carrier embedding
\[
m:Q\hookrightarrow P
\]
is split: there exists $r:P\to Q$ with $rm=\id_Q$.  In that case
\[
\tau=rp_r.
\]
\end{theorem}

\begin{proof}
If $\tau$ extends $h$, the maps $\tau:Q_r\to Q$ and $\id_Q:Q_*\to Q$ agree on $S$.  The pushout universal property gives a unique $r:P\to Q$ with
\[
rp_r=\tau,
\qquad
rm=\id_Q.
\]
Conversely, if $r$ is a retraction of $m$, then $\tau=rp_r$ satisfies, for $s\in S$,
\[
\tau(s)=rp_r(s)=rmh(s)=h(s).
\]
\end{proof}

\begin{theorem}[Forced closure as image containment]\label{thm:closurecontainment}
For a protected row, every row point is forced into the distinguished carrier if and only if
\[
p_r(Q)\subseteq m(Q),
\]
equivalently if and only if
\[
S^{\kappa}=Q.
\]
When this holds, $m^{-1}p_r:Q\to Q$ is the unique endomorphism extending $h$.
\end{theorem}

\begin{proof}
The equivalence with $S^\kappa=Q$ is Theorem~\ref{thm:kernelsat}.  If $p_r(Q)\subseteq m(Q)$, define $\tau=m^{-1}p_r$.  Since both $p_r$ and $m$ are homomorphisms and $m$ is injective,
$\tau$ is an endomorphism extending $h$. Every extension identifies $K$ and
therefore $\kappa$; if $x\equiv_\kappa s$ with $s\in S$, it must send $x$
to $h(s)$. Since $S^\kappa=Q$, these values determine the whole extension.
\end{proof}

Hence
\begin{equation}\label{eq:dictionary}
\boxed{
\begin{array}{rcl}
\text{protection}&\Longleftrightarrow&m\text{ monic},\\
\text{completion}&\Longleftrightarrow&m\text{ split monic},\\
\text{forced closure}&\Longleftrightarrow&p_r(Q)\subseteq m(Q).
\end{array}}
\end{equation}
Forced closure implies completion, and completion implies protection.
The converses fail. Example~\ref{ex:externalcompletion} gives a completion
with external row points. Example~\ref{ex:strictgap} below gives protection
without a completion on the protected carrier.

For several independent rows, strongly amalgamate all protected row factors over the same central $Q$.  A family of simultaneous endomorphism completions exists exactly when the single central embedding into this wide row amalgam splits.  Thus the local dictionary extends unchanged to the pure multirow fragment.

\subsection{Every repair lies above the forced quotient}

\begin{proposition}[Forced congruence is a universal lower bound]\label{prop:repairlower}
For every repair $\theta\in\mathcal R(E)$,
\[
\theta_E^A\subseteq\theta.
\]
\end{proposition}

\begin{proof}
A completion on $A/\theta$ is a model of the original ground theory in which the kernel of the named carrier map is exactly $\theta$.  Every provable carrier equality must therefore belong to $\theta$.
\end{proof}

Complete tables behave especially well: once the table is well defined on a quotient, protection already gives total endomorphisms, so the least repair is exactly the complete-table congruence of Corollary~\ref{cor:complete}.  Incomplete tables behave differently.

\subsection{No least repair}

\begin{example}[Two incomparable minimal repairs]\label{ex:noleastrepair}
Let $A=\{0,1,2,3\}$ have one unary operation
\[
f(0)=f(1)=f(2)=0,
\qquad
f(3)=1,
\]
and supply one row
\[
\alpha(b,1)=2.
\]
There is no completion on $A$.  Indeed, if $\tau$ were an endomorphism with $\tau(1)=2$, then $f(1)=0$ gives
\[
\tau(0)=\tau(f(1))=f(\tau(1))=f(2)=0,
\]
while $f(3)=1$ gives
\[
2=\tau(1)=\tau(f(3))=f(\tau(3)),
\]
impossible because $\im f=\{0,1\}$.

Nevertheless both
\[
\theta_a=\{0,2\}\mid\{1\}\mid\{3\}
\]
and
\[
\theta_b=\{0\}\mid\{1,2\}\mid\{3\}
\]
are repairs.  On $A/\theta_a$ the constant map to $[0,2]$ is an endomorphism satisfying the row; on $A/\theta_b$ the identity map satisfies it.  The congruences are incomparable and
\[
\theta_a\cap\theta_b=\Delta_A,
\]
which is not a repair. Each of $\theta_a,\theta_b$ merges only one pair,
so its only strictly finer equivalence relation is the diagonal. They are
therefore minimal repairs, and $\mathcal R(E)$ has no least element.
\end{example}

This failure is impossible for a complete table and is therefore a specifically incomplete-table phenomenon.

\subsection{Repair feasibility is not upward closed}

\begin{example}[A coarser quotient can destroy completion]\label{ex:nonmonotone}
Let $A=\{0,1,2,3\}$ have the constant unary operation $f(x)=0$ and supply
\[
\alpha(b,0)=1,
\qquad
\alpha(b,3)=2.
\]
The congruence
\[
\theta=\{0,1\}\mid\{2\}\mid\{3\}
\]
is a repair.  In the quotient, the only endomorphism constraint imposed by the constant operation is that $[0,1]$ map to itself; one may send $[3]$ to $[2]$.

Now coarsen to
\[
\phi=\{0,1,3\}\mid\{2\}.
\]
The two supplied source inputs $0$ and $3$ have become equal while their targets $1$ and $2$ remain distinct.  The projected row is not even a function, so $\phi$ is not a repair.  Total collapse is again a repair.  Thus $\mathcal R(E)$ is not an upper set of $\Con(A)$.
\end{example}

Consequently quotient repair cannot in general be found by a monotone search that simply coarsens until completion appears.

\subsection{The repair kernel and the cost of forbidding junk}

Define
\begin{equation}\label{eq:repairkernel}
\kappa_E^{\mathrm{NJ}}=
\bigcap_{\theta\in\mathcal R(E)}\theta.
\end{equation}
By Proposition~\ref{prop:repairlower},
\[
\theta_E^A\subseteq\kappa_E^{\mathrm{NJ}}.
\]
The inclusion can be strict.

\begin{example}[Strict repair-kernel gap]\label{ex:strictgap}
Let $A=\{0,1,2,3\}$ with
\[
f(0)=1,
\quad f(1)=1,
\quad f(2)=1,
\quad f(3)=2,
\]
and supply
\[
\alpha(b,0)=2,
\qquad
\alpha(b,1)=1,
\qquad
\alpha(b,2)=0.
\]
The original carrier is protected.  Indeed, adjoin one new element $x$, set $f(x)=0$, and define $\alpha(b,3)=x$. Give $\alpha(b,x)$ any value
to make the action total. Then all named-input compatibility equations hold while the four named carrier elements remain distinct.  Hence
\[
\theta_E^A=\Delta_A.
\]

There is no carrier-valued completion on $A$.  For any completion on a quotient, compatibility at input $3$ gives
\[
[0]=\tau([2])=\tau(f([3]))=f(\tau([3])).
\]
Thus $[0]$ must lie in the image of the quotient operation $f$, whose named image is represented by $[1]$ or $[2]$.  Hence $[0]=[1]$ or $[0]=[2]$.  In the first case the supplied rows at inputs $0$ and $1$ force $[2]=[1]=[0]$ as well.  Therefore every repair identifies $0$ and $2$.

The congruence
\[
\theta_*=\{0,2\}\mid\{1\}\mid\{3\}
\]
is itself a repair: the identity endomorphism of $A/\theta_*$ satisfies all projected rows.  Thus
\[
\boxed{
\theta_E^A=\Delta_A
<
\kappa_E^{\mathrm{NJ}}=\theta_*.
}
\]
\end{example}


The model-theoretic meaning of this gap can be stated without choosing any
completion procedure. Add the finite family of disjunctions
\begin{equation}\label{eq:nojunk}
E^{\mathrm{NJ}}=E\cup
\left\{\ \bigvee_{c\in A}\alpha(b,a)=c:
 b\in B,\ a\in A\ \right\}.
\end{equation}
This is a first-order theory of total algebras, not an equational theory.
Its additional condition concerns only the named action interface; it does
not require every element of a model to be named.

\begin{theorem}[Repairs as kernels of interface no-junk models]\label{thm:nojunk}
Writing $\ker_A(M)$ for the kernel of the named carrier map into $M$, one has
\begin{equation}\label{eq:repairmodels}
\mathcal R(E)=\{\ker_A(M):M\models E^{\mathrm{NJ}}\}.
\end{equation}
Consequently
\[
\kappa_E^{\mathrm{NJ}}
=\bigcap_{M\models E^{\mathrm{NJ}}}\ker_A(M),
\qquad
\theta_E^A=\bigcap_{M\models E}\ker_A(M)
\subseteq\kappa_E^{\mathrm{NJ}}.
\]
\end{theorem}

\begin{proof}
If $\theta$ is a repair, choose its endomorphism completions and interpret
the carrier sort by $A/\theta$ and the operator sort by $B$. This is a model
of $E^{\mathrm{NJ}}$ with named carrier kernel $\theta$.

Conversely, suppose $M\models E^{\mathrm{NJ}}$ and put
$\theta=\ker_A(M)$. For each $b$ and $[a]\in A/\theta$, choose a name $c$
with $\alpha^M(b^M,a^M)=c^M$, and set $\tau_b([a])=[c]$.
The choice of $c$ does not matter because all such names have the same
$\theta$-class. Nor does the representative $a$ matter, since equality of
its interpretations is preserved by the total operation $\alpha^M$.
Native ground diagrams and compatibility show that $\tau_b$ is an
endomorphism of $A/\theta$. The supplied cells hold by construction.
Thus $\theta$ is a repair. The intersection identities follow, the second
also being the usual model-theoretic characterization of ground equality.
\end{proof}

A pair in $\kappa_E^{\mathrm{NJ}}\setminus\theta_E^A$ can be kept distinct
only by allowing some named action value to be external. This is the
precise cost, in carrier distinctions, of requiring interface no-junk.
Example~\ref{ex:strictgap} shows that the cost can be nonzero even when the
original carrier is protected.

The no-junk condition is not generally preserved by direct products.
Each coordinate can choose a different carrier name for the same missing
cell, leaving its product value outside the diagonal named carrier image.
Thus a product can intersect carrier kernels while ceasing to satisfy
\eqref{eq:nojunk}. This explains why an intersection of repairs need not
itself be a repair, without asserting that repairs form a closure system.
